% =====================================================================
% Throughput Validation — Proposed Spark-Augmented Architecture
% =====================================================================
% Stand-alone subsection demonstrating that the proposed architecture
% (with Spark in the data path) sustains a useful production-grade
% throughput on the SIIT 12-core / 16 GB single-node hardware.
% =====================================================================

\subsection{Architecture Throughput Validation}
\label{sec:throughput}

To verify that the proposed architecture (Figure~\ref{fig:arch}),
including the Apache Spark preprocessor on the hot path, is
operationally usable rather than merely structurally complete, we
deployed the full multi-tenant stack on the
12-core / 16~GB Linux server and characterised its end-to-end
throughput. Synthetic flow records were produced by the load harness
(\texttt{loadtest/run\_experiment.py}, scenario~E1) at fixed offered
loads ranging from 50 to 1000~req/s, each sustained for 30~seconds
plus a 15-second drainage window. End-to-end latency was measured
from gateway POST to webhook arrival via per-record trace IDs.

\begin{figure}[t]
  \centering
  \includegraphics[width=\columnwidth]{figures/fig_throughput_sweep.pdf}
  \caption{Throughput sweep on the proposed Spark-augmented
  architecture. Bars (left axis) show end-to-end delivery rate; lines
  (right axis, log) show p50 and p99 latency. The pipeline sustains
  100\,\% delivery up to and including 200~req/s, after which Spark
  micro-batches saturate and delivery collapses.}
  \label{fig:throughput-sweep}
\end{figure}

\begin{table}[t]
\centering\small
\caption{Throughput sweep on the with-Spark architecture
(12-core / 16~GB single-node). Each row sustains the offered rate
for 30~s; latency is end-to-end gateway-to-webhook.}
\label{tbl:throughput}
\setlength{\tabcolsep}{4pt}
\begin{tabular}{rrrrl}
\toprule
\textbf{Rate} & \textbf{Delivered} & \textbf{p50 (ms)} & \textbf{p99 (ms)} & \textbf{Status} \\
\midrule
50      & 100\,\%   & 157   & 428      & comfortable     \\
100     & 100\,\%   & 144   & 256      & sustained        \\
200     & 100\,\%   & 3142  & 6340     & saturating      \\
500     & 23\,\%    & 2223  & $>$10\,min & collapsed   \\
1000    & 15\,\%    & $>$4\,min & $>$9\,min & overloaded \\
\bottomrule
\end{tabular}
\end{table}

\paragraph{Sustained throughput} The system delivered every flow
record through 100~req/s with a p99 latency of 256~ms, corresponding
to roughly $8.6$~million flows~per~tenant~per~day. At 100~req/s the
median latency (144~ms) is actually slightly \emph{lower} than at
50~req/s (157~ms): the Spark micro-batch executor accumulates more
records per trigger when the offered rate is higher, so the
per-record amortised cost falls. This is the operating regime that
production single-tenant deployments would target.

\paragraph{Burst capacity vs.\ saturation} At 200~req/s the pipeline
still preserved 100\,\% delivery but latency increased to multi-second
range (p50 = 3.1~s) because the detector consumer is unable to drain
each Spark micro-batch within one trigger interval. Beyond
approximately 300~req/s the gateway-Kafka-Spark-detector queue grows
without bound and delivery rates collapse to $<$25\,\%; we report
500 and 1000~req/s as overloaded conditions in
Table~\ref{tbl:throughput}. Higher per-tenant throughputs are
obtainable by horizontally scaling out the Spark executor and
Kafka partitions, which is the documented Spark cluster-mode path
and aligned with the cold-path scale-out justification in
Section~\ref{sec:retraining}.

\paragraph{Functional validation} Beyond the throughput numbers
themselves, this sweep confirms that the proposed architecture is
implementation-correct end-to-end: the Spark preprocessor (i)
subscribes to all per-tenant raw topics via a regex pattern,
(ii) applies the schema-normalisation, validation, and
feature-staging stages defined in Section~\ref{sec:proposed-system},
and (iii) republishes onto the per-tenant preprocessed topics that
the detector then consumes for downstream cascade scoring.

% =====================================================================
% End of section
% =====================================================================
